\input{../../../stock/en-tete_v6.tex}

\newcommand{\titrechapitre}{Convergence de suites de variables aléatoires}
\newcommand{\numerochapitre}{3}

\renewcommand{\headrulewidth}{0.4pt}
\renewcommand{\footrulewidth}{0.4pt}
\fancyfoot[L]{Notes de Raphaël JONTEF}
\fancyfoot[C]{\thepage}
\fancyfoot[R]{Enseirb-Matmeca}
\fancyhead[L]{Cours de Mathématiques}
\fancyhead[R]{\titrechapitre}

\begin{document}
\input{../../../stock/commands.tex}

% Page de titre custom
\setlength{\headheight}{15pt}
\begin{titlepage}
    \centering
    \vspace*{3cm}
    {\huge Chapitre \numerochapitre\par}
    {\Huge \bfseries\titrechapitre\par}
    \vspace{2cm}
    {\Large Notes rédigées par Raphaël JONTEF\par}
    \vspace{0.5cm}
    {\normalsize Cours enseigné par François DUFOUR\par}
    \vspace{4cm}
    {\small Enseirb-Matmeca\par}
    {\small filière informatique\par}
    \vfill
    {\large \today\par}
\end{titlepage}

% plan du cours
\tableofcontents
\newpage


% -----------Corps du document--------------

\section{Loi des grands nombres}

\begin{theoreme}{}{Loi des grands nombres}
    Soit $(X_n)_{n\in \N^*}$ une suite de variables aléatoires i.i.d. admettant un moment d'ordre 2. On note $\overline{X}_n = \frac{1}{n}\sum_{k=1}^n X_k$ sa moyenne empirique. Alors, pour tout $\varepsilon > 0$, on a~:
    $$\proba\Big( \vert \overline{X}_n - \esperance{X_1} \vert \geq \varepsilon\Big) \xrightarrow{n \to +\infty} 0$$
    Autrement dit, $\overline{X}_n$ converge en probabilité vers $\esperance{X_1}$.
    
\end{theoreme}

\begin{demonstration}
    Notons $\sigma^2$ la variance commune des $X_i$.    Par l'inégalité de Bienaymé-Tchebychev, on a~:
    $$\proba\Big( \vert \overline{X}_n - \esperance{X_1} \vert \geq \varepsilon\Big) \leq \frac{\mb{V}({\overline{X}_n})}{\varepsilon^2}$$
    Or par indépendance des variables aléatoires, on a~:
    $$\mb{V}({\overline{X}_n}) = \frac{1}{n^2}\sum_{i=1}^n\mb{V}({X_i}) = \frac{1}{n}\mb{V}({X_1})$$
    On en déduit que~:
    $$\proba\Big( \vert \overline{X}_n - \esperance{X_1} \vert \geq \varepsilon\Big) \leq \frac{\mb{V}({X_1})}{n\varepsilon^2}$$
    et le théorème s'en déduit.
\end{demonstration}

\begin{remarque}{}{convergence en probabilité}
    On peut alors écrire en terme de probabilité que pour presque tout $\omega \in \Omega$~:
    $$\frac{1}{n} \sum_{k=0}^nX_k(\omega) \xrightarrow{n \to +\infty} \esperance{X_1}$$
\end{remarque}

\begin{remarque}{}{applications}
    La loi des grands nombres permet d'approximer une espérance. Ceci présente des applications en finance.
\end{remarque}

\section{Différents modes de convergentes de suites de variables aléatoires}
Dans cette section, $k$ est un entier naturel non nul.
\begin{definition}{}{convergence en moyenne d'ordre $k$}
    Soit $(X_n)_{n \in \N}$ une suite de variables aléatoires admettant un moment d'ordre $k$ ($X_n^k$ est sommable ou intégrable, se note $x_n \in L^k$). On dit que $(X_n)_{n \in \N}$ converge vers $X \in L^k$ si 
    $$\lim_{n \to +\infty} \esperance{\vert X_n - X \vert^k} = 0$$
    On notera alors~:
    $$X_n \xrightarrow[\;n \to +\infty\;]{L^k} X$$
    
\end{definition}

\begin{definition}{}{convergence en probabilité}
    Soit $(X_n)_{n \in \N}$ une suite de variables aléatoires. On dit que $(X_n)_{n \in \N}$ converge en probabilité vers $X$ si pour tout $\varepsilon > 0$~:
    $$\lim_{n \to +\infty} \proba(\vert X_n - X \vert \geq \varepsilon) = 0$$
    On notera alors~:
        $$X_n \xrightarrow[\;n \to +\infty\;]{\proba} X$$
\end{definition}

\begin{definition}{}{convergence presque sûre}
    Soit $(X_n)_{n \in \N}$ une suite de variables aléatoires. On dit que $(X_n)_{n \in \N}$ converge presque sûrement vers $X$ si~:
    $$\proba\Big(\{\omega \in \Omega, \lim_{n \to +\infty} X_n(\omega) = X(\omega)\}\Big) = 1$$
    On notera alors~:
        $$X_n \xrightarrow[\;n \to +\infty\;]{\mr{p.s.}} X$$
    
\end{definition}
\newpage
On rappelle le théorème de convergence dominée pour les suites de fonctions définies sur un intervalle quelconque.

\begin{theoreme}{}{de convergence dominée pour une suite de fonctions définies sur un intervalle}
    Soit $(f_n)_{n \in \N} \in \mc{F}(I,\K)^\N$. Sous réserve des hypothèses suivantes,
    \begin{enumeratebf}
        \item $\forall n \in \N,\, f_n \in \mc{CM}(I,\K)$.
        \item $(f_n)_{n \in \N}$ converge simplement sur $I$ et sa limite y est continue par morceaux.
        \item \notion{hypothèse de domination} : il existe $\varphi \in \mc{CM}(I, \R_+)$ une fonction intégrable telle que pour tout $n \in \N$, $\abs{f_n} \leq \varphi$
    \end{enumeratebf}
    les fonctions $f_n$ sont intégrables, ainsi que leur limite simple et on peut échanger les symboles "$\lim$" et "$\int$"~:
    $$\lim_{n \to +\infty} \int_{I} f_n = \int_{I} \lim_{n \to +\infty} f_n $$
\end{theoreme}

\begin{theoreme}{}{de convergence dominée pour les variables aléatoires}
    Soit $(X_n)_{n \in \N}$ et $X$ des variables aléatoires telles que $X_n \xrightarrow[\;n \to +\infty\;]{\mr{p.s.}} X$. On suppose qu'il existe une variable aléatoire $Y$ intégrable telle que pour tout $n \in \N$, $\vert X_n \vert \leq Y$ presque sûrement. Alors $X$ est intégrable et~:
    $$\lim_{n \to +\infty} \esperance{X_n} = \esperance{X}$$
\end{theoreme}

\begin{proposition}{}{caractérisation de la convergence en probabilité par l'espérance}
    Soit $(X_n)_{n \in \N}$ une suite de variables aléatoires. Les propositions suivantes sont équivalentes~:
    \begin{enumeratebf}
        \item $X_n \xrightarrow[\;n \to +\infty\;]{\proba} X$
        \item $ \displaystyle \lim_{n \to +\infty} \esperance{\frac{\vert X_n - X\vert}{1 + \vert X_n - X\vert}} = 0$
    \end{enumeratebf}
\end{proposition}

\begin{remarque}{}{}
    Cette caractérisation est très puissante pour montrer une convergence en probabilité car elle repose sur l'intégration d'une fonction en $\frac{x}{1+x}$.
\end{remarque}

\begin{demonstration}
    \begin{itemize}
        \item $(1) \implies (2)$. On suppose que~:
        $$\forall \varepsilon > 0,\, \lim_{n \to +\infty} \probaset{\vert X_n - X \vert \geq \varepsilon} = 0$$
        Soit~:
        $$\forall \varepsilon > 0,\, \lim_{n \to +\infty} \esperance{\indicatrice_{\vert X_n - X \vert \geq \varepsilon}} = 0$$
        Puis pour $\varepsilon$ fixé~:
        \begin{align*}
            \esperance{\frac{\vert X_n - X\vert}{1 + \vert X_n - X\vert}} &= \esperance{\frac{\vert X_n - X\vert}{1 + \vert X_n - X\vert} \indicatrice_{\vert X_n - X \vert \geq \varepsilon}} + \esperance{\frac{\vert X_n - X\vert}{1 + \vert X_n - X\vert} \indicatrice_{\vert X_n - X \vert < \varepsilon}}\\
            &\leq  \esperance{\frac{\vert X_n - X\vert}{1 + \vert X_n - X\vert}} \\
            &= \esperance{\frac{\vert X_n - X\vert}{1 + \vert X_n - X\vert} \indicatrice_{\vert X_n - X \vert \geq \varepsilon}} + \frac{\varepsilon}{1+\varepsilon}
        \end{align*}

        \item 
    \end{itemize}


\end{demonstration}

\begin{proposition}{}{convergence presque sûre $\implies$ convergence en probabilité}
    Soit $(X_n)_{n \in \N}$ une suite de variables aléatoires convergeant presque sûrement vers une variable aléatoire $X$. Alors $(X_n)_{n \in \N}$ converge en probabilité vers $X$~:
    $$X_n \xrightarrow[\;n \to +\infty\;]{\mr{p.s}} X \quad\implies\quad X_n \xrightarrow[\;n \to +\infty\;]{\proba} X$$
\end{proposition}

\begin{demonstration}
    Puisque $X_n \xrightarrow[\;n \to +\infty\;]{\mr{p.s}} X$, alors 
\end{demonstration}

\begin{theoreme}{}{hiérarchie des convergences en moyenne}
    Soit $(X_n)_{n \in \N}$ une suite de variables aléatoires de $L^{k+1}$ convergent en moyenne vers $X$ à l'ordre $k+1$ pour $k \in \N^*$. Alors $(X_n)_{n \in \N}$ converge en moyenne vers $X$ à l'ordre $k$~:
    $$X_n \xrightarrow[\;n \to +\infty\;]{L^{k+1}} X \quad\implies\quad X_n \xrightarrow[\;n \to +\infty\;]{L^k} X$$
\end{theoreme}

\begin{theoreme}{}{convergence en moyenne d'ordre 1 $\implies$ convergence eeen probabilité}
    
\end{theoreme}

énorme trou.

\section{Loi des grands nombres}

\begin{theoreme}{}{loi faible des grands nombres}
    Soit $(X_n)_{n\in \N^*}$ une suite de variables aléatoires i.i.d. admettant un moment d'ordre 2. On note $\overline{X}_n = \frac{1}{n}\sum_{k=1}^n X_k$ sa \notion{moyenne empirique}. Alors, pour tout $\varepsilon > 0$, on a~:
    $$\proba\Big( \vert \overline{X}_n - \esperance{X_1} \vert \geq \varepsilon\Big) \xrightarrow{n \to +\infty} 0$$
    Autrement dit, $\overline{X}_n$ converge en probabilité vers $\esperance{X_1}$~:
    $$\overline{X}_n \xrightarrow[\;n \to +\infty\;]{\proba} \esperance{X_1}$$
    
\end{theoreme}

démo à copier


\begin{theoreme}{}{loi forte des grands nombres}
    Soit $(X_n)_{n\in \N^*}$ une suite de variables aléatoires i.i.d. admettant un moment d'ordre 1. On note $\overline{X}_n = \frac{1}{n}\sum_{k=1}^n X_k$ sa \notion{moyenne empirique}. Alors, pour presque tout $\omega \in \Omega$~:
    $$\frac{1}{n} \sum_{k=0}^nX_k(\omega) \xrightarrow{n \to +\infty} \esperance{X_1}$$
    Autrement dit, $\overline{X}_n$ converge presque sûrement vers $\esperance{X_1}$~:
    $$\overline{X}_n \xrightarrow[\;n \to +\infty\;]{\mr{p.s.}} \esperance{X_1}$$
\end{theoreme}

pas de démo

\section{Convergence faible et théorème de la limite centrale}
\subsection{Convergence faible}

\begin{definition}{}{convergence en loi}
    Soit $(X_n)_{n \in \N}$ et $X$ des variables aléatoires à valeurs dans $F$. On dit que $(X_n)_{n \in \N}$ converge en loi vers $X$ si pour toute fonction continue bornée $f : F \to G$, on a~:
    $$\lim_{n \to +\infty} \esperance{f(X_n)} = \esperance{f(X)}$$
    On notera alors~:
    $$ X_n \xrightarrow[\;n \to +\infty\;]{\mc{L}} X$$
    
\end{definition}

\subsection{Théorème de la limite centrale}
Soient $(X_n)_{n \in \N}$ une suite de variables aléatoires i.i.d. admettant un moment d'ordre 2. En notant $\sigma^2 = \mb{V}(X_1)$, on a~:
$$\frac{\sqrt{n}}{\sigma^2}\Big(\overline{X}_n - \esperance{X_1}\Big) \xrightarrow[\;n \to +\infty\;]{\mc{L}} Y$$
où $Y$ suit la loi normale centrée réduite $\mc{N}(0,1)$~:
\begin{itemize}
    \item centrée car son espérance et paramètre $\mu$ vaut 0
    \item réduite car sa variance vaut $\sigma^2 = 1$
\end{itemize}


\begin{remarque}{}{}
    D'un point de vue informatique, on peut montrer que générer toute variable aléatoire revient à savoir générer une variable aléatoire suivant une loi uniforme. On peut alors simuler n'importe quelle loi à partir d'une loi uniforme !
\end{remarque}
\section*{Questions de cours et remarques}

\subsection{Questions de cours}
\begin{itemize}
    \item Implications relatives aux convergences de suites de variables aléatoires et leur démonstration
    \item Loi faible des grands nombres et sa démonstration
    \item Loi forte des grands nombres
\end{itemize}

\end{document}